\documentclass[11pt, letterpaper]{article}
\usepackage[utf8]{inputenc}
\usepackage[spanish]{babel}
\usepackage{geometry}
\geometry{top=2.5cm, bottom=2.5cm, left=2.5cm, right=2.5cm}
\usepackage{amsmath}
\usepackage{hyperref}
\usepackage{xcolor}
\usepackage{listings}
\usepackage{enumitem}

\definecolor{codegray}{rgb}{0.5,0.5,0.5}
\definecolor{codelightgray}{rgb}{0.95,0.95,0.95}

\lstset{
    backgroundcolor=\color{codelightgray},
    commentstyle=\color{codegray},
    keywordstyle=\color{blue},
    numberstyle=\tiny\color{codegray},
    stringstyle=\color{purple},
    basicstyle=\ttfamily\footnotesize,
    breakatwhitespace=false,
    breaklines=true,
    captionpos=b,
    keepspaces=true,
    numbers=left,
    numbersep=5pt,
    showspaces=false,
    showstringspaces=false,
    showtabs=false,
    tabsize=4
}

\title{\textbf{Tarea 2: Análisis de partículas mediante estructuras dinámicas}}
\author{}
\date{}

\begin{document}

\begin{center}
    \LARGE \textbf{INF-73625 Programación Avanzada para Ciencias}\\
    \Large \textbf{Tarea 2: Análisis de partículas mediante estructuras dinámicas}
\end{center}

\vspace{0.5cm}

\noindent
\textbf{Profesores}
\begin{itemize}[noitemsep, topsep=0pt]
    \item Ricardo Salas Letelier (ricardo.salas@usm.cl)
    \item Wladimir Ormazabal (wladimir.ormazabal@usm.cl)
\end{itemize}

\vspace{0.3cm}
\noindent
\textbf{Ayudantes}
\begin{itemize}[noitemsep, topsep=0pt]
    \item Matías Gómez (mgomez@usm.cl)
\end{itemize}

\vspace{0.3cm}
\noindent
\textbf{Fecha de entrega:} 25 de septiembre del 2026

\section{Reglas}
La presente tarea debe hacerse de forma \textbf{individual}. No se permiten de ninguna manera trabajos en grupo.

Debe usarse el lenguaje de programación C++. Al evaluarla, la tarea será compilada usando el compilador \texttt{g++}, con la línea de comando \texttt{g++ archivo.cpp -o output -Wall}. Alternativamente, se aceptan variantes como MinGW.

Todas las estructuras de datos solicitadas deben ser implementadas completamente por cada estudiante. \textbf{No se permite usar la biblioteca stl}, así como ninguno de los containers y algoritmos definidos en ella. Está permitido usar bibliotecas estándar como \texttt{math.h}, \texttt{string}, \texttt{fstream}, \texttt{iostream}.

\section{Objetivos}
Entender y familiarizarse con la implementación y uso de estructuras de datos de tipo listas y árboles. Además, se fomentará el uso de buenas prácticas y el orden en la programación.

\section{Problemas a Resolver}
Entregue un archivo \texttt{.cpp} con una función \texttt{main} adecuada que permita probar su programa.

\subsection{Análisis de partículas mediante estructuras dinámicas}
En un experimento de física de partículas se registran distintos eventos en un detector. Dada la cantidad de información generada, se desea gestionar esta a través de estructuras de datos, donde se utilizará un Árbol Binario de Búsqueda (ABB) para almacenar los registros y una lista enlazada para las partículas que cumplan ciertas condiciones.

Cada partícula estará caracterizada por:
\begin{itemize}[noitemsep, topsep=0pt]
    \item \textbf{ID:} identificador único (string).
    \item \textbf{Energía:} (float).
    \item \textbf{Masa:} (float).
\end{itemize}
El ABB deberá estar ordenado de acuerdo con la energía de las partículas (se garantiza que no habrá dos partículas con la misma energía).

\subsubsection*{Formato de entrada}
Se entregará un archivo \texttt{Particulas.txt}. La primera línea contendrá un entero $N$ (cantidad de partículas). A continuación, se tendrán $N$ líneas con la información en el formato: \texttt{ID ENERGIA MASA}. Su programa deberá ser capaz de abrir este archivo y leer los datos para poblar las estructuras.

\textbf{Formato Particula.txt}

\begin{lstlisting}[language=bash]
8
P001 12.5 0.511
P002 7.8 1.20
P003 18.4 2.35
P004 4.2 0.105
P005 15.7 0.938
P006 22.1 4.18
P007 9.6 0.139
P008 13.3 0.494
\end{lstlisting}

\subsection{Construcción del Árbol Binario de Búsqueda}
Los registros deberán almacenarse en un ABB, utilizando la energía como criterio de orden. Implemente:
\begin{lstlisting}[language=C++]
void insertar(string id, float energia, float masa);
\end{lstlisting}
\begin{itemize}[noitemsep, topsep=0pt]
    \item \textbf{Entrada:} El ID de la partícula, su energía y su masa (leídos del archivo txt).
    
    \item \textbf{Resultado:} Modifica la estructura interna del árbol añadiendo un nuevo nodo.
\end{itemize}

Además, se deberá implementar un recorrido del árbol que permita visualizar la información almacenada en él. Implemente la función:
\begin{lstlisting}[language=C++]
void preOrden();
\end{lstlisting}
\begin{itemize}[noitemsep, topsep=0pt]
    \item \textbf{Resultado:} Imprime en consola un recorrido en \textbf{pre-orden} (raíz, izquierda, derecha) con la información completa de cada nodo, respetando el formato del archivo original.
\end{itemize}
\textbf{Ejemplo de salida esperada:}
\begin{lstlisting}[language=bash]
P001 12.5 0.511
P002 7.8 1.2
...
\end{lstlisting}

\subsection{Selección y construcción de la lista enlazada}
Para motivos del experimento, el usuario deberá definir dinámicamente un valor $U$ (umbral de energía). Se recorrerá el ABB seleccionando partículas con $\text{energía} \ge U$, incorporándolas a una lista enlazada con nodos independientes. Implemente:
\begin{lstlisting}[language=C++]
void insertarAlFinal(string id, float energia, float masa);
\end{lstlisting}
\begin{itemize}[noitemsep, topsep=0pt]
    \item \textbf{Entrada:} El ID, energía y masa de la partícula extraída del árbol.
    \item \textbf{Resultado:} Crea un nuevo nodo independiente y lo enlaza al final de la lista.
\end{itemize}

Una vez construida la lista, se deberá recorrer de manera secuencial e imprimir su contenido desde el inicio hasta el final. Implemente la función:
\begin{lstlisting}[language=C++]
void imprimirLista();
\end{lstlisting}
\begin{itemize}[noitemsep, topsep=0pt]
    \item \textbf{Resultado:} Imprime en consola \textbf{únicamente el ID} de las partículas separadas por flechas.
\end{itemize}
\textbf{Ejemplo de salida esperada:}
\begin{lstlisting}[language=bash]
P006 -> P003 -> P005 -> NULL
\end{lstlisting}

\subsection{Eliminación de partículas de la lista}
Permita eliminar una partícula por ID (por ejemplo, para descartar lecturas ruidosas). Implemente:
\begin{lstlisting}[language=C++]
void eliminar(string idEliminar);
\end{lstlisting}
\begin{itemize}[noitemsep, topsep=0pt]
    \item \textbf{Entrada:} El identificador único de la partícula a descartar (ej. \texttt{"P001"}).
    \item \textbf{Resultado:} Modifica la estructura de la lista liberando la memoria del nodo. Debe considerar los casos: inicio, medio, final, lista vacía y elemento no encontrado.
\end{itemize}

\subsection{Inversión Magnética de Partículas}
Durante el análisis de las partículas seleccionadas en la lista, el detector fue expuesto accidentalmente a un campo magnético inverso severo[cite: 3]. Como consecuencia, el orden cronológico de las partículas en la lista enlazada quedó exactamente opuesto[cite: 3]. Implemente la función:
\begin{lstlisting}[language=C++]
void InvertirLista();
\end{lstlisting}
\begin{itemize}[noitemsep, topsep=0pt]
    \item \textbf{Resultado:} Reasigna los punteros internos de los nodos para que el orden de la lista se invierta físicamente.
\end{itemize}


\textbf{Ejemplo de funcionamiento:}
Si la lista original era:
\begin{lstlisting}[language=bash]
P006 -> P003 -> P005 -> NULL
\end{lstlisting}

Después de llamar a esta función, una llamada a \texttt{imprimirLista()} deberá mostrar:
\begin{lstlisting}[language=bash]
P005 -> P003 -> P006 -> NULL
\end{lstlisting}

\section{Entrega de la Tarea}
Entregue la tarea enviando un archivo comprimido tarea1-nombre-apellido.zip (reemplazando por 
nombre y apellido según corresponda) a la página aula.usm.cl del curso, a más tardar el día 28 de 
septiembre 2026, a las 23:59 (Chile Continental), el cual contenga:
\begin{itemize}[noitemsep, topsep=0pt]
    \item Los archivos con los códigos fuentes necesarios para el funcionamiento de la tarea. ¡Los 
archivos deben compilar!

    
    \item \textbf{README.txt:} Instrucciones de compilación en caso de ser necesarias, y la forma de 
compilación que usó (debe ser alguna de las indicadas en los tutoriales entregados en 
Aula).
\end{itemize}


\section{Restricciones y consideraciones}
Por cada día de atraso se descontarán 10 puntos en la nota. Las tareas que no compilen serán calificadas con nota 0. Por cada Warning, -5 puntos. Si se detecta COPIA la nota automáticamente será 0.

\end{document}